\documentclass[11pt]{article}

\usepackage[margin=0.9in]{geometry}
\usepackage{amsmath,amssymb}
\usepackage{booktabs}
\usepackage{enumitem}
\usepackage{listings}
\usepackage{xcolor}
\usepackage{hyperref}
\usepackage{tikz}

\definecolor{backcolour}{rgb}{0.95,0.95,0.92}
\definecolor{codeblue}{rgb}{0.05,0.15,0.55}
\definecolor{codegreen}{rgb}{0.0,0.45,0.0}

\lstset{
    backgroundcolor=\color{backcolour},
    basicstyle=\ttfamily\small,
    keywordstyle=\color{codeblue}\bfseries,
    commentstyle=\color{codegreen},
    stringstyle=\color{purple},
    breaklines=true,
    frame=single,
    showstringspaces=false,
    tabsize=4
}

\title{CPSC 250 \\ Classes Program Example: The GVDie Craps Game}
\author{Edward Brash}
\date{}

\begin{document}

\maketitle

\section{Program Goal}

In this example, we write a program to play an automated dice game using a provided \texttt{GVDie} class.

The player rolls two dice. Depending on the result, the player may:

\begin{itemize}
    \item win one credit,
    \item lose one credit,
    \item or set a goal for future rolls.
\end{itemize}

A round ends when the player either wins or loses a credit. The entire game ends when the player reaches zero credits.

\section{The Provided \texttt{GVDie} Class}

The provided die class represents a single six-sided die.

A simplified version is:

\begin{lstlisting}[language=Python]
import random

class GVDie:

    def __init__(self):
        self.value = None

    def roll(self):
        self.value = random.randint(1, 6)

    def get_value(self):
        return self.value

    def compare_to(self, d):
        return self.value == d.value
\end{lstlisting}

\section{The Meaning of \texttt{None}}

The constructor initializes:

\begin{lstlisting}[language=Python]
self.value = None
\end{lstlisting}

The value \texttt{None} means that the die does not yet have a meaningful value.

This is useful because a die should not really have a value until it has been rolled.

\section{Importing and Using the Die Class}

If the class is stored in a file called \texttt{gv\_die.py}, then the main program can import it:

\begin{lstlisting}[language=Python]
import random
from gv_die import GVDie
\end{lstlisting}

Then we can create die objects:

\begin{lstlisting}[language=Python]
die1 = GVDie()
die2 = GVDie()
\end{lstlisting}

Each die object has its own internal value.

\section{Random Seeds}

The program may ask the user for a random seed:

\begin{lstlisting}[language=Python]
seed = int(input("Enter seed: "))
random.seed(seed)
\end{lstlisting}

A seed defines the particular sequence of random numbers.

This is useful for testing because it makes program behavior repeatable.

\section{Credits}

The user also enters the number of credits:

\begin{lstlisting}[language=Python]
credits = int(input("Enter starting credits: "))
\end{lstlisting}

These credits can be interpreted as the player's number of lives.

The game continues until:

\begin{lstlisting}[language=Python]
credits == 0
\end{lstlisting}

\section{Basic Game Setup}

The notes outline the beginning of the program as follows:

\begin{enumerate}
    \item Create two \texttt{GVDie} objects.
    \item Initialize the number of rounds.
    \item Initialize the goal value to \texttt{-1}.
    \item Begin the main game loop.
\end{enumerate}

\begin{lstlisting}[language=Python]
die1 = GVDie()
die2 = GVDie()

rounds = 0
goal = -1

while credits > 0:
    ...
\end{lstlisting}

\section{How a Round Works}

A round begins with an initial roll of both dice.

\begin{lstlisting}[language=Python]
die1.roll()
die2.roll()

roll = die1.get_value() + die2.get_value()
\end{lstlisting}

The total determines what happens next.

\section{Initial Roll Outcomes}

For the initial roll:

\begin{itemize}
    \item rolling 7 or 11 wins one credit,
    \item rolling 3 or 12 loses one credit,
    \item any other value becomes the goal.
\end{itemize}

\begin{lstlisting}[language=Python]
if roll == 7 or roll == 11:
    credits = credits + 1

elif roll == 3 or roll == 12:
    credits = credits - 1

else:
    goal = roll
\end{lstlisting}

\section{Goal Rolls}

If the initial roll sets a goal, the player continues rolling.

On later rolls:

\begin{itemize}
    \item rolling 7 loses one credit,
    \item rolling the goal wins one credit,
    \item any other roll means roll again.
\end{itemize}

\begin{lstlisting}[language=Python]
while goal != -1:

    die1.roll()
    die2.roll()

    roll = die1.get_value() + die2.get_value()

    if roll == 7:
        credits = credits - 1
        goal = -1

    elif roll == goal:
        credits = credits + 1
        goal = -1
\end{lstlisting}

\section{Why This Is a Class Example}

This program is useful because it shows objects in action.

Instead of treating dice as raw integers, we use die objects that know how to:

\begin{itemize}
    \item roll themselves,
    \item store their current value,
    \item return their value,
    \item compare themselves to another die.
\end{itemize}

The class bundles together both data and behavior.

\section{Complete Program Skeleton}

\begin{lstlisting}[language=Python]
import random
from gv_die import GVDie

seed = int(input("Enter seed: "))
random.seed(seed)

credits = int(input("Enter starting credits: "))

die1 = GVDie()
die2 = GVDie()

rounds = 0

while credits > 0:

    rounds = rounds + 1
    goal = -1

    die1.roll()
    die2.roll()

    roll = die1.get_value() + die2.get_value()

    if roll == 7 or roll == 11:
        credits = credits + 1

    elif roll == 3 or roll == 12:
        credits = credits - 1

    else:
        goal = roll

    while goal != -1:

        die1.roll()
        die2.roll()

        roll = die1.get_value() + die2.get_value()

        if roll == 7:
            credits = credits - 1
            goal = -1

        elif roll == goal:
            credits = credits + 1
            goal = -1

print("Game over")
print("Rounds played:", rounds)
\end{lstlisting}

\section{Key Takeaways}

\begin{itemize}
    \item Objects combine state and behavior.
    \item A die object stores its current value.
    \item Methods such as \texttt{roll()} and \texttt{get\_value()} define how users interact with the object.
    \item A random seed makes simulations repeatable.
    \item Larger programs become manageable when broken into loops, branches, and objects.
\end{itemize}

\end{document}
